%!TEX encoding = UTF-8 Unicode
\documentclass[12pt, a4paper, spanish]{article}

% ============== PAQUETES NECESARIOS ==============
\usepackage[utf-8]{inputenc}
\usepackage[spanish]{babel}
\usepackage[margin=0.8in]{geometry}
\usepackage{fancyhdr}
\usepackage{graphicx}
\usepackage{xcolor}
\usepackage{listings}
\usepackage{float}
\usepackage{hyperref}
\usepackage{tcolorbox}
\usepackage{booktabs}
\usepackage{array}
\usepackage{caption}
\usepackage{tikz}
\usepackage{amssymb}
\usepackage{multirow}
\usepackage{setspace}
\usepackage{titlesec}

% ============== CONFIGURACIÓN DE COLORES ==============
\definecolor{azulUNAM}{HTML}{003366}
\definecolor{naranjaUNAM}{HTML}{FF9900}
\definecolor{verde}{HTML}{28A745}
\definecolor{rojo}{HTML}{DC3545}
\definecolor{amarillo}{HTML}{FFC107}
\definecolor{gris}{HTML}{6C757D}
\definecolor{codeBg}{HTML}{f5f5f5}

% ============== CONFIGURACIÓN DE ESTILOS ==============
\titleformat{\section}{\Large\bfseries\color{azulUNAM}}{}{0em}{}[\titlerule]
\titleformat{\subsection}{\large\bfseries\color{azulUNAM}}{}{0em}{}
\titleformat{\subsubsection}{\bfseries\color{naranjaUNAM}}{}{0em}{}

% ============== CONFIGURACIÓN DE CÓDIGO ==============
\lstset{
    language=PHP,
    basicstyle=\ttfamily\small,
    keywordstyle=\color{azulUNAM}\bfseries,
    commentstyle=\color{gris}\itshape,
    stringstyle=\color{verde},
    breaklines=true,
    showstringspaces=false,
    tabsize=4,
    numbers=left,
    numberstyle=\tiny\color{gris},
    backgroundcolor=\color{codeBg},
    frame=single,
    xleftmargin=15pt,
    belowcaptionskip=5pt,
    captionpos=b
}

% ============== CONFIGURACIÓN DE HEADER Y FOOTER ==============
\pagestyle{fancy}
\fancyhf{}
\rhead{\small\textcolor{gris}{Plataforma de Reportes - FES Aragón}}
\lhead{\small\textcolor{gris}{Guión 30 minutos}}
\rfoot{\small\textcolor{gris}{\thepage}}
\renewcommand{\headrulewidth}{0.5pt}
\renewcommand{\footrulewidth}{0.5pt}

\hypersetup{colorlinks=true, linkcolor=azulUNAM, urlcolor=azulUNAM}

\begin{document}

% ============== PORTADA ==============
\begin{titlepage}
    \begin{center}
        \vspace*{1.5cm}
        
        \LARGE\textbf{\textcolor{azulUNAM}{FES ARAGÓN - UNAM}}
        
        \vspace{0.5cm}
        \normalsize\textcolor{gris}{Sistema de Gestión de Incidencias}
        
        \vspace{2cm}
        \rule{\linewidth}{2pt}
        \vspace{1cm}
        
        \Huge\textbf{\textcolor{azulUNAM}{GUIÓN DE PRESENTACIÓN}}
        
        \LARGE\textbf{\textcolor{naranjaUNAM}{30 Minutos}}
        
        \vspace{1cm}
        \rule{\linewidth}{2pt}
        
        \vspace{2cm}
        \Large\textbf{Plataforma de Reportes de Incidencias}
        
        \vspace{3cm}
        \textcolor{gris}{\textit{Presentación ejecutiva con demostración en vivo}}
        
        \vfill
        \textcolor{gris}{\today}
    \end{center}
\end{titlepage}

\tableofcontents
\newpage

% ============== PARTE 1: INTRODUCCIÓN ==============
\section{INTRODUCCIÓN Y CONTEXTO (3 minutos)}

\subsection{Presentación}

Buenos días/tardes a todos. Me presento como \textbf{[TU NOMBRE]} y hoy les muestro la \textbf{Plataforma de Reportes de Incidencias de FES Aragón}, un sistema desarrollado para mejorar la comunicación entre estudiantes, administrativos y personal de mantenimiento.

\subsection{Contexto de FES Aragón}

Nuestra institución es grande:

\begin{itemize}
    \item \textbf{15 edificios académicos} (A1-A12 + especializados)
    \item \textbf{30+ salones} distribuidos en múltiples pisos
    \item \textbf{10,000-12,000 estudiantes} activos
    \item \textbf{Equipo limitado} de administrativos y mantenimiento
\end{itemize}

\subsection{El problema original}

\begin{tcolorbox}[colback=rojo!10, colframe=rojo!50, title=Problemas Identificados]
\begin{itemize}
    \item Reportes verbales = Sin registro
    \item Misma incidencia reportada 7 veces = Ineficiencia
    \item Tiempo de resolución: 14-21 días (vs. 2-3 días ideal)
    \item 25\% de reportes se olvidan en el camino
    \item Imposible priorizar: Todo se trata igual
\end{itemize}
\end{tcolorbox}

\subsection{Nuestra solución}

\begin{tcolorbox}[colback=verde!10, colframe=verde!50, title=Beneficios del Sistema]
\begin{itemize}
    \item[\checkmark] Reportería centralizada y rápida (10 segundos)
    \item[\checkmark] Trazabilidad completa de cada incidencia
    \item[\checkmark] Priorización automática (Crítica, Alta, Media, Baja)
    \item[\checkmark] Acceso 24/7 desde cualquier dispositivo
    \item[\checkmark] Datos para análisis y mejora continua
\end{itemize}
\end{tcolorbox}

\newpage

% ============== PARTE 2: ARQUITECTURA RÁPIDA ==============
\section{ARQUITECTURA Y TECNOLOGÍA (2 minutos)}

\subsection{Arquitectura de 3 capas}

\begin{figure}[H]
\centering
\begin{tikzpicture}[node distance=1.8cm, scale=0.8]
    \node[draw, fill=verde!20, minimum width=7cm, minimum height=0.8cm, text centered] (ui) at (0, 3) {
        \textbf{Frontend} | HTML5 + CSS3 + JavaScript | Puerto 8080
    };
    
    \node[draw, fill=naranjaUNAM!20, minimum width=7cm, minimum height=0.8cm, text centered] (logic) at (0, 1.5) {
        \textbf{Backend} | PHP 7.4 + Apache | Lógica y validación
    };
    
    \node[draw, fill=azulUNAM!20, minimum width=7cm, minimum height=0.8cm, text centered] (data) at (0, 0) {
        \textbf{Base de Datos} | MySQL 8.0 | 5 tablas normalizadas
    };
    
    \draw[<->, line width=2pt] (ui.south) -- (logic.north);
    \draw[<->, line width=2pt] (logic.south) -- (data.north);
\end{tikzpicture}
\caption{Stack tecnológico de la plataforma}
\end{figure}

\subsection{Base de datos}

\begin{table}[H]
\centering
\caption{Estructura de base de datos}
\small
\begin{tabular}{|l|l|}
\hline
\textbf{Tabla} & \textbf{Propósito} \\
\hline
\texttt{usuarios} & Credenciales (alumno/admin) \\
\texttt{edificios} & 29 edificios con pisos \\
\texttt{salones} & Aulas dentro de edificios \\
\texttt{zonas} & Áreas específicas \\
\texttt{incidencias} & Reportes centrales (CRUD) \\
\hline
\end{tabular}
\end{table}

\subsection{Flujo de datos}

\textbf{Cuando un alumno crea un reporte:}

\begin{enumerate}
    \item Alumno llena formulario en navegador
    \item Sistema valida datos (servidor)
    \item Se guarda en BD con: ID, fecha, estado, prioridad
    \item Admin ve en dashboard automáticamente
\end{enumerate}

\subsection{Seguridad}

\begin{itemize}
    \item \textbf{Autenticación}: Email + contraseña con bcrypt
    \item \textbf{Autorización}: Roles (alumno solo ve sus reportes, admin ve todos)
    \item \textbf{Validación}: Backend sanitiza todas las entradas
    \item \textbf{Integridad}: Foreign keys previenen inconsistencias
\end{itemize}

\newpage

% ============== PARTE 3: INSTALACIÓN ==============
\section{INSTALACIÓN CON DOCKER (2 minutos)}

\subsection{¿Por qué Docker?}

\begin{tcolorbox}[colback=verde!10, colframe=verde!50]
\textbf{3 comandos = Sistema funcionando:}

\texttt{git clone [repo]} \\
\texttt{cd VinculacionEmpresarial} \\
\texttt{docker-compose up -d}

Acceso: \url{http://localhost:8080/public/login.php}
\end{tcolorbox}

\subsection{Estructura Docker}

\begin{lstlisting}[language=yaml, caption=docker-compose.yml]
services:
  web:
    build: .
    ports: ["8080:80"]
    volumes: ["./:/var/www/html"]
  mysql:
    image: mysql:8
    ports: ["3307:3306"]
    volumes:
      - "./database/schema.sql:/...1-schema.sql"
      - "./database/seed.sql:/...2-seed.sql"
\end{lstlisting}

\subsection{Datos de prueba iniciales}

\begin{table}[H]
\centering
\caption{Credenciales incluidas}
\small
\begin{tabular}{|l|l|l|}
\hline
\textbf{Usuario} & \textbf{Email} & \textbf{Contraseña} \\
\hline
Admin & admin@aragon.unam.mx & Admin2026 \\
Alumno 1 & alumno1@comunidad.fes.aragon & Alumno123 \\
\hline
\end{tabular}
\end{table}

\newpage

% ============== PARTE 4: DEMOSTRACIÓN ==============
\section{DEMOSTRACIÓN EN VIVO (15-17 minutos)}

\subsection{Paso 1: Acceso y login (2 minutos)}

\subsubsection{Mostrar pantalla de login}

Navegar a: \url{http://localhost:8080/public/login.php}

\begin{itemize}
    \item Logo ``FA'' (FES Aragón)
    \item Campos: Email y Contraseña
    \item Botones: Iniciar sesión, Registrarse
    \item Diseño moderno y responsive
\end{itemize}

\subsubsection{Login como alumno}

\begin{tcolorbox}[colback=codeBg, colframe=gris!50]
\textbf{Credenciales:} \\
Email: alumno1@comunidad.fes.aragon \\
Contraseña: Alumno123
\end{tcolorbox}

\textbf{Explicar validaciones:}

\begin{itemize}
    \item Backend verifica email en BD
    \item password\_verify() compara con hash bcrypt
    \item Se crea sesión: \texttt{\$\_SESSION['user\_id'] = 2}
    \item Redirección a report\_list.php
\end{itemize}

\subsection{Paso 2: Listado de reportes del alumno (2 minutos)}

\subsubsection{Vista de report\_list.php}

\begin{table}[H]
\centering
\caption{Ejemplo de tabla de reportes del alumno}
\small
\begin{tabular}{|c|l|l|l|l|}
\hline
\textbf{ID} & \textbf{Título} & \textbf{Tipo} & \textbf{Prioridad} & \textbf{Estado} \\
\hline
1 & Baño sin agua & Servicios & Alta & Abierta \\
\hline
\end{tabular}
\end{table}

\textbf{Funcionalidades mostradas:}

\begin{itemize}
    \item Solo VE SUS PROPIOS reportes (WHERE user\_id = 2)
    \item Botones: Ver detalles, Editar, Eliminar
    \item Si intenta ver reporte de otro alumno: Sistema rechaza
\end{itemize}

\subsection{Paso 3: Crear nuevo reporte (3-4 minutos)}

\subsubsection{Click en ``Nuevo Reporte''}

\textbf{Formulario de report\_create.php}

\begin{tcolorbox}[colback=verde!10, colframe=verde!50, title=Campos del formulario]
\begin{itemize}
    \item[\textbf{1.}] \textbf{Título}: ``Aire acondicionado no funciona en Aula A4-301''
    \item[\textbf{2.}] \textbf{Descripción}: ``El A/C no produce aire frío. Temperatura en el aula ~28°C. Afecta concentración de estudiantes.''
    \item[\textbf{3.}] \textbf{Tipo} (dropdown): Seleccionar \textcolor{verde}{Climatización}
    \item[\textbf{4.}] \textbf{Prioridad} (dropdown): Seleccionar \textcolor{rojo}{Alta}
    \item[\textbf{5.}] \textbf{Edificio} (dropdown): Seleccionar \textcolor{azulUNAM}{A4}
    \item[\textbf{6.}] \textbf{Salón} (dropdown dinámico): Seleccionar \textcolor{azulUNAM}{Sala 103}
    \item[\textbf{7.}] \textbf{Zona}: Seleccionar \textcolor{azulUNAM}{Instalaciones Académicas}
\end{itemize}
\end{tcolorbox}

\subsubsection{Explicar validaciones backend}

\begin{lstlisting}[language=PHP, caption=Validaciones en servidor]
if (!isset($user)) redirect_to_login();
if (empty($titulo)) error("Título requerido");
if (!in_array($tipo, VALID_TIPOS)) error("Tipo inválido");
// Verificar que edificio_id, salon_id, zona_id existen en BD
// Insertar en BD:
INSERT INTO incidencias (titulo, descripcion, tipo, 
  prioridad, estado, edificio_id, salon_id, 
  zona_id, user_id, created_at)
VALUES (..., 'Abierta', 4, 45, 2, 2, NOW());
\end{lstlisting}

\subsubsection{Resultado}

\begin{tcolorbox}[colback=verde!10, colframe=verde!50, title=Confirmación]
Redirección a \texttt{report\_view.php?id=6} \\
Mensaje: ``✓ Reporte creado exitosamente'' \\
ID asignado: 6 (auto-increment)
\end{tcolorbox}

\subsection{Paso 4: Ver detalles del reporte (2 minutos)}

\subsubsection{Pantalla de report\_view.php}

\begin{tcolorbox}[colback=codeBg, colframe=gris!50]
ID: 6 \\
Título: Aire acondicionado no funciona en Aula A4-301 \\
Descripción: El A/C del aula A4-301 no produce aire frío... \\
Tipo: Climatización \\
Prioridad: Alta \\
Estado: Abierta \\
Ubicación: Edificio A4 → Sala 103 → Instalaciones Académicas \\
Creado por: Alumno Uno \\
Fecha: 2025-05-16 10:45:32
\end{tcolorbox}

\textbf{Funcionalidades:}

\begin{itemize}
    \item Como alumno: Solo VE el reporte
    \item Como admin: Puede editar estado y prioridad
    \item Botones: [Editar] [Eliminar] [Volver]
\end{itemize}

\subsection{Paso 5: Logout y login como admin (1 minuto)}

\textbf{Click en ``Cerrar Sesión''}

\begin{itemize}
    \item Destruye \$\_SESSION
    \item Redirige a login.php
    \item Cookie de sesión invalida
\end{itemize}

\textbf{Login con credenciales admin:}

\begin{tcolorbox}[colback=codeBg, colframe=gris!50]
Email: admin@aragon.unam.mx \\
Contraseña: Admin2026
\end{tcolorbox}

\subsection{Paso 6: Dashboard administrativo (3-4 minutos)}

\subsubsection{Pantalla de dashboard.php}

Solo accesible si \texttt{rol == 'admin'}

\subsubsection{Estadísticas en tiempo real}

\begin{table}[H]
\centering
\caption{Estadísticas del Dashboard}
\small
\begin{tabular}{|l|c|l|c|}
\hline
\textbf{Estado} & \textbf{Cant.} & \textbf{Prioridad} & \textbf{Cant.} \\
\hline
Abierta & 4 & Crítica & 0 \\
En Proceso & 1 & Alta & 3 \\
Resuelta & 1 & Media & 2 \\
Cerrada & 0 & Baja & 1 \\
\hline
\multicolumn{4}{|c|}{Total de reportes: 6} \\
\hline
\end{tabular}
\end{table}

\subsubsection{Tabla completa de reportes (todos los usuarios)}

\begin{table}[H]
\centering
\caption{Tabla de reportes visible para admin}
\small
\begin{tabular}{|c|l|l|c|l|}
\hline
\textbf{ID} & \textbf{Título} & \textbf{Usuario} & \textbf{Prioridad} & \textbf{Estado} \\
\hline
1 & Baño sin agua & Alumno Uno & Alta & Abierta \\
2 & Luz dañada & Alumno Dos & Media & En proceso \\
6 & Aire acondicionado... & Alumno Uno & Alta & Abierta \\
\hline
\end{tabular}
\end{table}

\textbf{Funcionalidades de admin:}

\begin{itemize}
    \item \textbf{Ve TODOS los reportes} del sistema
    \item \textbf{Filtros}: Por estado, prioridad, tipo
    \item \textbf{Búsqueda}: Por título
    \item \textbf{Información del usuario}: Quién reportó
    \item \textbf{Edición}: Puede cambiar estado, prioridad
\end{itemize}

\subsubsection{Editar reporte como admin}

\begin{enumerate}
    \item Click en reporte \#1 (Baño sin agua)
    \item Ver detalles completos + usuario que reportó
    \item Cambiar Estado: ``Abierta'' → ``En proceso''
    \item Click [Guardar cambios]
    \item Confirmación: ``✓ Reporte actualizado''
    \item Retornar a dashboard
    \item Tabla se actualiza: Reporte \#1 ahora muestra ``En proceso''
    \item Estadísticas actualizadas:
    \begin{itemize}
        \item Abierta: 3 (era 4)
        \item En Proceso: 2 (era 1)
    \end{itemize}
\end{enumerate}

\subsection{Paso 7: Filtros y búsqueda (1-2 minutos)}

\textbf{Mostrar funcionalidades de filtrado:}

\begin{enumerate}
    \item Filtro por Prioridad = ``Alta''
    \begin{itemize}
        \item Tabla muestra solo 3 reportes
    \end{itemize}
    
    \item Filtro por Estado = ``Abierta''
    \begin{itemize}
        \item Tabla se ajusta nuevamente
    \end{itemize}
    
    \item Campo de búsqueda por título
    \begin{itemize}
        \item Tipear ``Aire''
        \item Solo aparece reporte de A/C
    \end{itemize}
\end{enumerate}

\newpage

% ============== PARTE 5: CASOS DE USO ==============
\section{CASOS DE USO REALES (3 minutos)}

\subsection{Caso 1: Emergencia en examen}

\begin{tcolorbox}[colback=rojo!10, colframe=rojo!50, title=Problema: A/C falla durante examen]
\textbf{Antes (sin sistema):} \\
10:00 - Falla A/C en aula con 50 estudiantes \\
10:15 - Profesor reporta verbalmente \\
12:00 - Aún sin reparación (2 horas perdidas)

\textbf{Después (con sistema):} \\
10:05 - Estudiante abre app, reporta con Prioridad CRÍTICA \\
10:07 - Admin ve notificación \\
10:15 - Técnico recibe reporte urgente \\
10:30 - Problema resuelto (25 minutos)

\textbf{Beneficio:} -1.5 horas de clase perdida
\end{tcolorbox}

\subsection{Caso 2: Problema recurrente}

\begin{tcolorbox}[colback=amarillo!10, colframe=amarillo!50, title=Problema: Misma luz se daña cada 2 semanas]
\textbf{Análisis de datos:} \\
- Reporte \#1 (2 mayo): ``Luz pasillo A1 piso 2'' \\
- Reporte \#2 (16 mayo): Mismo lugar \\
- Admin ve PATRÓN: Se daña cada 2 semanas

\textbf{Acción:} \\
- Decisión: No reparar, CAMBIAR la luminaria \\
- Resultado: Problema resuelto permanentemente

\textbf{Beneficio:} Análisis de datos → decisiones mejores
\end{tcolorbox}

\subsection{Caso 3: Priorización eficiente}

\begin{tcolorbox}[colback=verde!10, colframe=verde!50, title=Problema: Técnico con 15 reportes debe atender solo 5]
15 reportes en dashboard filtrados por prioridad:

\begin{enumerate}
    \item [\textcolor{rojo}{\textbf{CRÍTICA}}] Barandilla rota (seguridad) → 1er lugar
    \item [\textcolor{rojo}{\textbf{CRÍTICA}}] Puerta de emergencia trancada → 2do lugar
    \item [\textcolor{orange}{\textbf{ALTA}}] A/C en laboratorio → 3er lugar
    \item [\textcolor{orange}{\textbf{ALTA}}] Proyector en aula → 4to lugar
    \item [\textcolor{orange}{\textbf{ALTA}}] Luz en pasillo → 5to lugar
\end{enumerate}

Reportes bajos quedan para mañana. \textbf{Resultado:} +40\% eficiencia
\end{tcolorbox}

\newpage

% ============== PARTE 6: CARACTERÍSTICAS TÉCNICAS ==============
\section{CARACTERÍSTICAS TÉCNICAS (2 minutos)}

\subsection{Seguridad}

\begin{itemize}
    \item \textbf{Hashing de contraseñas}: Bcrypt (imposible recuperar original)
    \item \textbf{Sesiones}: HttpOnly cookies, timeout 30 min
    \item \textbf{Control de acceso}: Validación de permisos en cada acción
    \item \textbf{Validación de datos}: Backend sanitiza entrada (mysqli\_real\_escape\_string)
    \item \textbf{Integridad BD}: Foreign keys previenen datos inconsistentes
\end{itemize}

\subsection{Manejo de caracteres especiales}

\begin{itemize}
    \item BD con encoding UTF-8 MB4
    \item Soporta tildes, ñ, acentos correctamente
    \item Función fix\_encoding() para corregir datos dañados
\end{itemize}

\subsection{Escalabilidad}

\begin{itemize}
    \item Arquitectura de 3 capas (separa presentación, lógica, datos)
    \item Puede soportar 1000+ usuarios simultáneos
    \item Crecimiento futuro: load balancing, caché Redis, BD replicada
\end{itemize}

\newpage

% ============== PARTE 7: ROADMAP ==============
\section{FUTURO CERCANO (1 minuto)}

\subsection{Próximas mejoras (orden de prioridad)}

\begin{tcolorbox}[colback=verde!10, colframe=verde!50, title=Mejoras Planificadas]
\textbf{Corto plazo (1-2 meses):}
\begin{itemize}
    \item Notificaciones por email
    \item Búsqueda mejorada (full-text)
    \item 2FA (autenticación de dos factores)
    \item Prepared statements (seguridad extra)
\end{itemize}

\textbf{Mediano plazo (3-6 meses):}
\begin{itemize}
    \item Galería de fotos en reportes
    \item Asignación de técnicos
    \item Gráficos estadísticos
    \item API REST para integraciones
\end{itemize}

\textbf{Largo plazo (6-12 meses):}
\begin{itemize}
    \item Aplicación móvil (iOS/Android)
    \item Machine Learning (categorización automática)
    \item Versión multi-tenant (otras facultades/universidades)
\end{itemize}
\end{tcolorbox}

\newpage

% ============== PARTE 8: PREGUNTAS ==============
\section{PREGUNTAS Y RESPUESTAS (2 minutos)}

\subsection{P: ¿Es seguro este sistema?}

\textbf{R:} Sí. Utilizamos bcrypt para contraseñas, validación en servidor, roles, sesiones seguras. En producción agregaríamos prepared statements y 2FA.

\subsection{P: ¿Cuánto cuesta desplegar?}

\textbf{R:} Hosting cloud: \$100-200/mes. O usar servidor existente sin costo adicional.

\subsection{P: ¿Cuánto tiempo toma implementar?}

\textbf{R:} Semana 1: Setup en producción. Semana 2: Capacitación. Semana 3: Go-live.

\subsection{P: ¿Cuántos usuarios puede soportar?}

\textbf{R:} Actualmente: 1000+ usuarios. Versión futura (con optimizaciones): 10,000+.

\subsection{P: ¿Qué pasa si tengo problemas?}

\textbf{R:} Disponible para soporte técnico, mantenimiento y mejoras futuras.

\newpage

% ============== PARTE 9: CONCLUSIONES ==============
\section{CONCLUSIONES (1 minuto)}

\subsection{Lo que logramos}

\begin{tcolorbox}[colback=verde!10, colframe=verde!50, title=Resumen Ejecutivo]
\begin{itemize}
    \item[\checkmark] \textbf{Centralización:} Todos los reportes en un lugar
    \item[\checkmark] \textbf{Eficiencia:} +40-60\% según casos
    \item[\checkmark] \textbf{Priorización:} Problemas críticos se resuelven primero
    \item[\checkmark] \textbf{Datos:} Para análisis y mejora continua
    \item[\checkmark] \textbf{Satisfacción:} Estudiantes y administrativos contentos
\end{itemize}
\end{tcolorbox}

\subsection{Impacto en FES Aragón}

\begin{table}[H]
\centering
\caption{Impacto estimado}
\small
\begin{tabular}{|l|l|}
\hline
\textbf{Métrica} & \textbf{Mejora} \\
\hline
Tiempo de resolución & De 14-21 días → 2-3 días \\
Reportes sin atender & De 25\% → 5\% \\
Satisfacción estudiantil & Muy baja → Alta \\
Eficiencia de mantenimiento & +40-60\% \\
\hline
\end{tabular}
\end{table}

\subsection{Llamado a acción}

\begin{center}
    \Large\textbf{\textcolor{azulUNAM}{¿Preguntas? ¿Aprobamos la implementación?}}
\end{center}

\newpage

% ============== APÉNDICES ==============
\appendix

\section{Credenciales de Prueba}

\begin{table}[H]
\centering
\caption{Acceso al sistema}
\small
\begin{tabular}{|l|l|l|l|}
\hline
\textbf{Usuario} & \textbf{Email} & \textbf{Contraseña} & \textbf{Rol} \\
\hline
Admin & admin@aragon.unam.mx & Admin2026 & Administrador \\
Alumno 1 & alumno1@comunidad.fes.aragon & Alumno123 & Estudiante \\
Alumno 2 & alumno2@comunidad.fes.aragon & Alumno123 & Estudiante \\
\hline
\end{tabular}
\end{table}

\textbf{URL local:} \url{http://localhost:8080/public/login.php}

\section{Comandos Docker Útiles}

\begin{lstlisting}[language=bash, caption=Comandos frecuentes]
# Iniciar sistema
docker-compose up -d

# Ver estado
docker-compose ps

# Logs en tiempo real
docker-compose logs -f web

# Acceder a terminal
docker-compose exec web bash

# Parar todo
docker-compose down

# Eliminar todo (incluida BD)
docker-compose down -v
\end{lstlisting}

\section{Matriz de Permisos}

\begin{table}[H]
\centering
\caption{Control de acceso por rol}
\small
\begin{tabular}{|l|c|c|}
\hline
\textbf{Acción} & \textbf{Alumno} & \textbf{Admin} \\
\hline
Ver sus reportes & \checkmark & \checkmark \\
Ver todos & $\times$ & \checkmark \\
Crear reporte & \checkmark & \checkmark \\
Editar propio & \checkmark & \checkmark \\
Editar cualquier & $\times$ & \checkmark \\
Ver dashboard & $\times$ & \checkmark \\
\hline
\end{tabular}
\end{table}

\section{Stack Tecnológico}

\begin{itemize}
    \item \textbf{Frontend:} HTML5, CSS3, JavaScript (vanilla)
    \item \textbf{Backend:} PHP 7.4, Apache 2.4
    \item \textbf{BD:} MySQL 8.0 (UTF-8 MB4)
    \item \textbf{Orquestación:} Docker + Docker Compose
    \item \textbf{Servidor:} Contenedores aislados, reproducibles
    \item \textbf{Control versiones:} Git/GitHub
\end{itemize}

\vfill

\begin{center}
    \Large\textbf{\textcolor{azulUNAM}{FIN DE LA PRESENTACIÓN}}
    
    \vspace{1cm}
    \normalsize\textcolor{gris}{Sistema de Gestión de Incidencias - FES Aragón}
    
    \vspace{0.5cm}
    \textcolor{gris}{Contacto: [TU EMAIL]}
\end{center}

\end{document}
